\documentclass[11pt,a4paper]{article}
\usepackage[margin=1in]{geometry}
\usepackage{hyperref}
\usepackage[T1]{fontenc}
\usepackage[utf8]{inputenc}
\usepackage[english]{babel}
\usepackage{enumitem}
\usepackage{longtable}
\usepackage{array}
\usepackage{booktabs}
\usepackage{xcolor}

% Hyperlink setup
\hypersetup{
    colorlinks=true,
    linkcolor=blue,
    urlcolor=blue,
    citecolor=blue
}

% Remove page numbers
\pagestyle{plain}

% Adjust spacing
\setlength{\parindent}{0pt}
\setlength{\parskip}{0.5em}

% Compact metadata block for each publication
\newcommand{\pubmeta}[3]{%
  \par\smallskip\textit{\small\textbf{Status:} #1; \textbf{Indexing:} #2; \textbf{Defended provisions claimed:} #3.}\par}

\begin{document}

% Header
\begin{center}
{\LARGE \textbf{List of Published Scientific Works on the Dissertation Topic}}\\[0.4em]
{\large Roger Nick Anaedevha (A23 -- 501)}\\[0.3em]
PhD candidate, Department of Cybernetics No.\,22, \\ Institute of Intelligent Cybernetic Systems (IICS)\\
National Research Nuclear University ``MEPhI'', Moscow, Russia\\[0.3em]
\textit{Dissertation title:} ``Development of Adversarial Resilient Models Based on Artificial Intelligence in Hybrid Intrusion Detection and Prevention Systems for Network Security''\\[0.3em]
ORCID: \href{https://orcid.org/0000-0003-3649-1561}{0000-0003-3649-1561} ${\cdot}$
\href{https://scholar.google.com/citations?user=w6ryqsEAAAAJ&hl=en}{Google Scholar} ${\cdot}$
\href{https://github.com/rogerpanel/}{Github.com/rogerpanel/}\\
\href{mailto:ar006@campus.mephi.ru}{ar006@campus.mephi.ru} ${\cdot}$ \href{mailto:rogernickanaedevha@gmail.com}{rogernickanaedevha@gmail.com}
\end{center}

\vspace{-0.4cm}

%%---------------------------------------------------------------------
\section*{Notation: Defended Provisions}
%%---------------------------------------------------------------------

In accordance with the section ``Main scientific results submitted for defence'' of the dissertation:

\begin{itemize}[leftmargin=2em, itemsep=0.1em]
  \item \textbf{Provision No.\,1} --- Models M1 (CT-TGNN, SDE-TGNN) for certified dynamics of the learnable detector of the Hybrid~IDPS (invariant P1).
  \item \textbf{Provision No.\,2} --- Algorithm M2 (FedLLM-API) for privacy-preserving distributed training of the Hybrid~IDPS (invariant P4).
  \item \textbf{Provision No.\,3} --- Architecture M3 (TripleE-TGNN) for multi-level data representation in the Hybrid~IDPS (invariant P6).
  \item \textbf{Provision No.\,4} --- Architecture M4 (MambaShield) for efficient inference of the Hybrid~IDPS (invariant P3).
  \item \textbf{Provision No.\,5} --- Models M5--M6 (stochastic transformer and UC-HGP) for drift-aware uncertainty-calibrated adaptation of the Hybrid~IDPS (invariants P5, P6).
  \item \textbf{Provision No.\,6} --- Unified Hybrid~IDPS certification framework (M7) that integrates models M1--M6 through Stackelberg game-theoretic equilibria and is validated on network traffic (Chapter~5) and on UAV onboard systems (Chapter~6).
\end{itemize}

\vspace{-0.2cm}

\textit{Note:} this list includes \textbf{only published} (in-print) journal articles and conference proceedings on the dissertation topic. Preprints, manuscripts under review, and works accepted for publication but not yet in print at the time this list was compiled are deliberately excluded.

\vspace{-0.2cm}

%%---------------------------------------------------------------------
\section*{1. Published Journal Articles}
%%---------------------------------------------------------------------

\begin{enumerate}[leftmargin=*, itemsep=0.3em]

\item \textbf{Anaedevha, R.N., Trofimov, A.G., Borodachev, Y.V. (2025).} ``MambaShield: Selective State-Space Models for Poisoning-Resilient Sequential Modeling.'' \textit{Expert Systems with Applications}, Elsevier. DOI: \href{https://doi.org/10.1016/j.eswa.2026.131175}{10.1016/j.eswa.2026.131175}. ISSN 0957-4174.
\pubmeta{published}{Scopus Q1, WoS Core (SCIE), VAK (K1), RSCI; JCR IF\,$\approx$\,7.5}{No.\,4 (M4 -- MambaShield)}

\item \textbf{Anaedevha, R.N., Trofimov, A.G., Borodachev, Y.V. (2026).} ``Uncertainty-calibrated hierarchical Gaussian processes for intrusion detection with multi-scale temporal modeling.'' \textit{Neurocomputing}, Elsevier, p.\,133105. DOI: \href{https://doi.org/10.1016/j.neucom.2026.133105}{10.1016/j.neucom.2026.133105}. ISSN 0925-2312.
\pubmeta{published}{Scopus Q1, WoS Core (SCIE), VAK (K1), RSCI; JCR IF\,$\approx$\,5.5}{No.\,5 (M5--M6 -- stochastic transformer and UC-HGP)}

\item \textbf{Anaedevha, R.N., Trofimov, A.G., Borodachev, Y.V. (2026).} ``Temporal Adaptive Neural Ordinary Differential Equations with Deep Spatio-Temporal Point Processes for Real-Time Network Intrusion Detection.'' \textit{Complex and Intelligent Systems}, Springer. DOI: \href{https://doi.org/10.1007/s40747-026-02291-7}{10.1007/s40747-026-02291-7}. ISSN 2199-4536.
\pubmeta{published}{Scopus Q1, WoS Core (SCIE), VAK (K1), RSCI; JCR IF\,$\approx$\,5.0}{No.\,1 (M1 -- CT-TGNN, SDE-TGNN)}

\item \textbf{Anaedevha, R.N., Trofimov, A.G. (2024).} ``Improved Robust Adversarial Model against Evasion Attacks on Intrusion Detection Systems.'' \textit{Optical Memory and Neural Networks (Information Optics)}, Vol.\,33, Suppl.\,3, pp.\,S414--S423, Springer. DOI: \href{https://doi.org/10.3103/S1060992X24700681}{10.3103/S1060992X24700681}. ISSN 1060-992X.
\pubmeta{published}{Scopus, VAK (K2), RSCI}{No.\,6 (M7 -- unified certification framework, foundational adversarial-robustness formulation)}

\end{enumerate}

\vspace{-0.3cm}
%%---------------------------------------------------------------------
\section*{2. Published Conference Proceedings}
%%---------------------------------------------------------------------

\begin{enumerate}[leftmargin=*, itemsep=0.3em, resume]

\item \textbf{Anaedevha, R.N., Trofimov, A.G. (2025).} ``Stochastic Multimodal Transformer with Uncertainty Quantification for Robust Network Intrusion Detection.'' In: \textit{Advances in Neural Computation, Machine Learning, and Cognitive Research IX. NEUROINFORMATICS 2025.} Studies in Computational Intelligence, vol.\,1241, Springer, Cham. DOI: \href{https://doi.org/10.1007/978-3-032-07690-8_35}{10.1007/978-3-032-07690-8\_35}. ISBN 978-3-032-07689-2.
\pubmeta{published}{Scopus (Springer SCI), RSCI}{No.\,5 (M5 -- stochastic multimodal transformer with uncertainty quantification)}

\item \textbf{Anaedevha, R.N., Trofimov, A.G. (2025).} ``Integrated Deep Q-Networks and Actor-Critic Models for Enhanced Poisoning Attack Detection in Network Intrusion Detection Systems.'' In: \textit{2025 Conference of Young Researchers in Electrical and Electronic Engineering (ElCon)}, ETU-LETI, Saint-Petersburg, Russia. IEEE Xplore, vol.\,33, suppl.\,3, pp.\,556--567.
\pubmeta{published}{IEEE Xplore, Scopus, RSCI}{No.\,6 (M7 -- game-theoretic and reinforcement-learning response component)}

\item \textbf{Anaedevha, R.N. (2024).} ``Application of Machine Learning Models for Device Identification in Wireless Network Traffic.'' In: \textit{2024 Conference of Young Researchers in Electrical and Electronic Engineering (ElCon)}, Saint-Petersburg, Russia. IEEE Xplore, pp.\,104--110. DOI: \href{https://doi.org/10.1109/ElCon61730.2024.10468413}{10.1109/ElCon61730.2024.10468413}.
\pubmeta{published}{IEEE Xplore, Scopus, RSCI}{No.\,1 (M1 -- foundational feature-extraction problem statement for the NIDPS detector)}

\end{enumerate}

\vspace{-0.2cm}

%%---------------------------------------------------------------------
\section*{3. Summary Counts by Type and Level}
%%---------------------------------------------------------------------

\textbf{Total published scientific works on the dissertation topic: 7}, of which:

\begin{longtable}{p{0.62\textwidth} >{\centering\arraybackslash}p{0.10\textwidth} >{\centering\arraybackslash}p{0.20\textwidth}}
\toprule
\textbf{Category} & \textbf{Count} & \textbf{Share of 7} \\
\midrule
\endhead

\multicolumn{3}{l}{\textit{By type of edition}} \\
\midrule
Journal articles & 4 & 57.1\% \\
Conference proceedings (full text, indexed in Scopus / IEEE Xplore) & 3 & 42.9\% \\
\midrule

\multicolumn{3}{l}{\textit{By indexing (all 7 published works)}} \\
\midrule
Scopus Q1, WoS Core (SCIE), VAK (K1) --- journal articles & 3 & 42.9\% \\
Scopus, VAK (K2), RSCI --- journal article & 1 & 14.3\% \\
Scopus (Springer SCI), RSCI --- conference series & 1 & 14.3\% \\
IEEE Xplore (Scopus), RSCI --- conference proceedings & 2 & 28.6\% \\
\midrule

\multicolumn{3}{l}{\textit{By authorship}} \\
\midrule
First author & 7 & 100.0\% \\
\midrule

\multicolumn{3}{l}{\textit{By Higher Attestation Commission (VAK) requirements (minimum 3 for engineering)}} \\
\midrule
Journal articles in editions on the VAK list and in Scopus/WoS & 4 & --- \\
Of those, in editions of VAK category K1 and Scopus Q1 & 3 & --- \\
\bottomrule
\end{longtable}

\textbf{Compliance with the Regulation on the Award of Academic Degrees (Government of the Russian Federation Decree No.\,842 of 24.09.2013, paragraphs 11--13):} the requirement on the minimum number of publications in peer-reviewed scientific editions (not less than 3 for engineering sciences) is \textbf{fulfilled and exceeded}: 4 published journal articles, of which 3 are in Q1 Scopus / WoS Core / VAK~K1 editions.

\vspace{-0.2cm}

%%---------------------------------------------------------------------
\section*{4. Cross-Reference: Publications vs. Defended Provisions}
%%---------------------------------------------------------------------

\begin{longtable}{>{\raggedright\arraybackslash}p{0.22\textwidth} >{\raggedright\arraybackslash}p{0.50\textwidth} >{\centering\arraybackslash}p{0.22\textwidth}}
\toprule
\textbf{Provision} & \textbf{Published works} & \textbf{Additional confirmation} \\
\midrule
\endhead

\textbf{No.\,1} (M1: CT-TGNN, SDE-TGNN) &
[3]~Anaedevha, R.N. \textit{et al.}, \textit{Complex \& Intelligent Systems}, Springer, 2026 (Q1) \newline
[7]~Anaedevha, R.N., ElCon 2024, IEEE Xplore (foundational) &
Implementation act No.\,1 \linebreak (NII\,IKT LLC) \\
\midrule

\textbf{No.\,2} (M2: FedLLM-API) &
\emph{(no published works on this provision at the time of compiling this list; supported by implementation acts --- see Section~5)} &
Implementation acts No.\,1 and No.\,2 \linebreak (NII\,IKT LLC, Hexagon LLC) \\
\midrule

\textbf{No.\,3} (M3: TripleE-TGNN) &
\emph{(no published works on this provision at the time of compiling this list; supported by an implementation act --- see Section~5)} &
Implementation act No.\,3 \linebreak (AI Center, NRNU MEPhI, Moscow) \\
\midrule

\textbf{No.\,4} (M4: MambaShield) &
[1]~Anaedevha, R.N. \textit{et al.}, \textit{Expert Systems with Applications}, Elsevier, 2025 (Q1) &
Implementation act No.\,2 \linebreak (Hexagon LLC) \\
\midrule

\textbf{No.\,5} (M5--M6: stochastic transformer, UC-HGP) &
[2]~Anaedevha, R.N. \textit{et al.}, \textit{Neurocomputing}, Elsevier, 2026 (Q1) \newline
[5]~Anaedevha, R.N., Trofimov, A.G., NEUROINFORMATICS 2025, Springer SCI &
Implementation act No.\,3 \linebreak (AI Center, NRNU MEPhI, Moscow) \\
\midrule

\textbf{No.\,6} (M7: unified certification framework, game-theoretic) &
[4]~Anaedevha, R.N., Trofimov, A.G., OMNN, Springer, 2024 (foundational adversarial robustness) \newline
[6]~Anaedevha, R.N., Trofimov, A.G., ElCon 2025, IEEE Xplore (RL response) &
Implementation act No.\,1 \linebreak (NII\,IKT LLC) \\
\bottomrule
\end{longtable}

\textit{Conclusion:} four of the six defended provisions (No.\,1, No.\,4, No.\,5, No.\,6) are directly supported by published works; provisions No.\,2 (M2 FedLLM-API) and No.\,3 (M3 TripleE-TGNN) are, at the time of compiling this list, supported empirically~--- by implementation acts at NII\,IKT LLC, Hexagon LLC, and the AI Center of NRNU MEPhI (see Section~5). When both forms of corroboration are combined (publications and implementation acts), coverage of all six defended provisions reaches \textbf{100\,\%}.

\vspace{-0.2cm}

%%---------------------------------------------------------------------
\section*{5. Implementation and Usage Acts}
%%---------------------------------------------------------------------

\textbf{Three implementation acts} have been received that confirm the practical applicability of the developed adversarial resilient models based on artificial intelligence in hybrid intrusion detection and prevention systems for network security:

\begin{enumerate}[leftmargin=*, itemsep=0.4em]

\item \textbf{Implementation act from NII\,IKT LLC} (Research Institute of Information and Communication Technologies), Moscow. \\
\textit{Subject of implementation:} application of the developed AI-based Hybrid~IDPS framework (models M1, M2, M7) to network-security protection of critical-infrastructure facilities. \\
\textit{Confirmed performance:} 95.7\,\% detection accuracy on streaming traffic from critical-infrastructure facilities in a distributed environment. \\
\textit{Provisions confirmed:} No.\,1, No.\,2, No.\,6.

\item \textbf{Implementation act from Hexagon LLC}, Moscow. \\
\textit{Subject of implementation:} integration of federated learning (M2) and efficient edge inference (M4) into commercial Hybrid~IDPS deployments with support for distributed-environment operation and edge-computing security. \\
\textit{Confirmed performance:} 96.8\,\% detection accuracy and a 40\,\% reduction in computational cost relative to the reference architecture. \\
\textit{Provisions confirmed:} No.\,2, No.\,4.

\item \textbf{Implementation act from the Artificial Intelligence Center of the National Research Nuclear University ``MEPhI''}, Moscow. \\
\textit{Subject of implementation:} application of the developed AI-based Hybrid~IDPS framework (models M1, M3, M6) to the protection of UAV onboard systems within the three-tier ``onboard~-- droneport~-- cloud'' architecture (corresponds to Chapter~6 of the dissertation). \\
\textit{Confirmed performance:} confirmation of domain independence of the developed AI-based framework via transfer from the network-traffic sector (NIDPS) to the adjacent UAV onboard-system sector. \\
\textit{Provisions confirmed:} No.\,1, No.\,3, No.\,5, No.\,6.

\end{enumerate}

\textit{Outcome:} taken together, the three implementation acts confirm \emph{all} six defended provisions.

\vspace{-0.2cm}

%%---------------------------------------------------------------------
\section*{6. Software Products and Intellectual-Property Registration}
%%---------------------------------------------------------------------

\begin{enumerate}[leftmargin=*, itemsep=0.3em]

\item \textbf{Software platform ``RobustIDPS.ai''}~--- a unified open-source software platform that integrates 17 neural-network models into 10 functional groups and 51 pages of a Progressive Web Application user interface; implements all models M1--M7 as components of the unified Hybrid~IDPS system, including a four-layer large-language-model defence framework (94.5\,\% detection across 517 attack scenarios) and the Multi-Agent PQC-IDS module. \\
\textit{Source code deposit:} Zenodo, DOI \href{https://doi.org/10.5281/zenodo.19129512}{10.5281/zenodo.19129512}. \\
\textit{Provisions confirmed:} No.\,1--No.\,6 (all six, jointly within the unified system).

\item \textbf{Software-for-computers (SW) registration applications with Rospatent} are planned after dissertation defence for the key components of ``RobustIDPS.ai''~--- CT-TGNN, FedLLM-API, TripleE-TGNN, MambaShield, and the unified M7 certification framework (a total of 5 applications).

\item \textbf{Invention patents:} none have been filed at the time of preparation of this list; applications are planned following commercial validation at Hexagon LLC for the Multi-Agent PQC-IDS and FedLLM-API components.

\end{enumerate}

\vspace{0.4cm}
\hrule
\vspace{0.2cm}

\small
\textit{Last updated: June 2026.}\\
\textit{Abbreviations:} ``K1'', ``K2''~--- categories of peer-reviewed scientific editions on the list of the Higher Attestation Commission (VAK) under the Ministry of Science and Higher Education of the Russian Federation; ``Q1'', ``Q2''~--- journal quartiles in the Scopus / Web of Science databases at the time of manuscript submission; ``IF''~--- two-year impact factor from Journal Citation Reports (JCR); ``RSCI''~--- Russian Science Citation Index (RINTs).

\textit{Complete publication records and citation metrics are maintained on the author's Google Scholar and ORCID profiles.}

\end{document}
